\subsection*{Models}
\label{sec:axiom-systems-models}

\begin{toolkitbox}{Models Toolkit --- Quick Reference}
\small
\begin{tabular}{l l l}
\toprule
\textbf{Concept} & \textbf{Role} & \textbf{Detail} \\
\midrule
Interpretation & assigns meanings to symbols & \hyperref[def:axiom-systems-interpretation]{Definition} \\
Structure & domain with interpreted symbols & \hyperref[def:axiom-systems-structure]{Definition} \\
Model & structure satisfying the theory & \hyperref[def:axiom-systems-model]{Definition} \\
Satisfaction & truth of a statement in a structure & \hyperref[def:axiom-systems-satisfaction]{Definition} \\
Countermodel & model where a statement fails & \hyperref[def:axiom-systems-countermodel]{Definition} \\
\bottomrule
\end{tabular}
\end{toolkitbox}

\begin{tcolorbox}[colback=structurebox, colframe=structureborder, arc=2pt,
  left=8pt, right=8pt, top=6pt, bottom=6pt,
  title={\small\textbf{Model Presentation: Arithmetic as a Model}},
  fonttitle=\small\bfseries]
Reuse the arithmetic signature and arithmetic-style axioms from the previous
sections. A model begins by interpreting the symbols of that signature in an
actual mathematical structure.

\[
\begin{array}{rcl}
\text{Domain} & = & \text{the natural numbers},\\
0 & \mapsto & \text{the number zero},\\
S & \mapsto & \text{the successor operation},\\
+ & \mapsto & \text{addition}.
\end{array}
\]

In this structure, the statement
\[
  \forall x\;(x+0=x)
\]
holds because adding zero to any natural number gives back that same natural
number.

The formal symbol $0$ is only a symbol until it is interpreted as the number
zero. The formal symbol $+$ is only a symbol until it is interpreted as
addition. The symbols acquire meaning only after interpretation.
\end{tcolorbox}

\begin{remark*}[Structural remarks]
\textbf{Syntax vs meaning.}
Signatures, axioms, and theories are formal objects. They belong to syntax:
they specify vocabulary, adopted assertions, and consequences. Models provide
semantic meaning by assigning mathematical content to the symbols and by
determining which statements hold.

\textbf{One theory, many models.}
The same theory may have multiple models. Euclidean geometry may be realized
in different coordinate systems. The group axioms have many models, including
different concrete groups. The theory fixes the assertions; a model is one
mathematical realization in which those assertions hold.

\textbf{Countermodels.}
A statement fails to follow from a theory if there is a model of the theory in
which that statement is false. Countermodels are therefore evidence that a
claim is not forced by the axioms.

\textbf{Forward roadmap.}
The next sections use models to explain consistency and independence:
\[
\begin{array}{c}
\text{Models}\\
\Downarrow\\
\text{Consistency}\\
\Downarrow\\
\text{Independence}
\end{array}
\]
\end{remark*}

\begin{definition}[Interpretation]
\label{def:axiom-systems-interpretation}
An \emph{interpretation} is an assignment of meanings to the non-logical
symbols of a signature.
\end{definition}

\begin{remark*}[Interpretation]
An interpretation turns formal vocabulary into mathematical content. It tells
what each constant, function symbol, and relation symbol means in a particular
setting.
\end{remark*}

\begin{dependencies}
\begin{itemize}
  \item \hyperref[def:signature]{Signature}
\end{itemize}
\end{dependencies}

\begin{definition}[Structure]
\label{def:axiom-systems-structure}
A \emph{structure} is a domain together with interpretations of the symbols in
a signature.
\end{definition}

\begin{remark*}[Interpretation]
A structure is the environment in which formal symbols receive meaning. The
domain supplies the objects, and the interpretation explains how the symbols
refer to objects, operations, and relations on that domain.
\end{remark*}

\begin{dependencies}
\begin{itemize}
  \item \hyperref[def:signature]{Signature}
  \item \hyperref[def:axiom-systems-interpretation]{Interpretation}
\end{itemize}
\end{dependencies}

\begin{definition}[Model]
\label{def:axiom-systems-model}
A \emph{model} is a structure in which all axioms, equivalently all statements
of the theory, are satisfied.
\end{definition}

\begin{remark*}[Interpretation]
A model is a structure that makes the theory true. It is not just an
interpretation of the vocabulary; it is an interpretation in which the adopted
assertions and their consequences hold.
\end{remark*}

\begin{dependencies}
\begin{itemize}
  \item \hyperref[def:axiom-systems-structure]{Structure}
  \item \hyperref[def:axiom-systems-theory]{Theory}
\end{itemize}
\end{dependencies}

\begin{definition}[Satisfaction]
\label{def:axiom-systems-satisfaction}
A statement is \emph{satisfied} in a structure if it evaluates as true under
the given interpretation.
\end{definition}

\begin{remark*}[Interpretation]
Satisfaction is the semantic test for a statement. Once the symbols have been
interpreted in a structure, the statement can be checked as true or false in
that setting.
\end{remark*}

\begin{dependencies}
\begin{itemize}
  \item \hyperref[def:axiom-systems-structure]{Structure}
\end{itemize}
\end{dependencies}

\begin{definition}[Countermodel]
\label{def:axiom-systems-countermodel}
A \emph{countermodel} is a model of a theory in which a particular statement
fails.
\end{definition}

\begin{remark*}[Interpretation]
A countermodel shows that the theory does not force the statement. It keeps
the axioms true while making the proposed statement false.
\end{remark*}

\begin{dependencies}
\begin{itemize}
  \item \hyperref[def:axiom-systems-model]{Model}
  \item \hyperref[def:axiom-systems-satisfaction]{Satisfaction}
\end{itemize}
\end{dependencies}
